\section{Project Context}
\label{sec:context}

\subsection{Client and Problem}

The client is \ClientName{}, who acts as the Product Owner.
In Guayaquil, people and businesses that need to store goods or inventory find storage space through informal deals.
They get no guarantee about the safety of the place or its legal compliance.
At the same time, owners of warehouses with idle space have no formal channel to offer and manage it.
Leodega is a two-sided marketplace that connects both groups: managers publish warehouses, an administrator verifies them, and clients search for them and book them.
The agreed scope is recorded in the Project Scope Document, signed on 2026-10-04 (Section~\ref{sec:client}).

\subsection{System Scope}

The system has one backend and two clients: the Leodega web app and Leodeguita, its mobile companion.
Both clients use the same accounts, data and business rules.
The MVP covers:
\begin{itemize}
  \item warehouse registration with a fire department permit, plus editing and deletion;
  \item administrator review of each warehouse before it is published;
  \item catalog search with filters and a map, and booking with date conflict detection;
  \item organizations with members, invitations, a personal/organization context selector and bookings made for the organization;
  \item later sprints: chat, incident reports, ratings, external inventory integration and an administrator dashboard.
\end{itemize}
Payment is simulated.
The client first named the Paganini gateway, but another course team builds it and it is not ready, so that integration is postponed to a later phase.

\subsection{User Roles}

\begin{itemize}
  \item \textbf{Administrator}: approves or rejects warehouses and their fire permits, blocks or reactivates accounts, and resolves reports.
  \item \textbf{Storage Manager} (\textit{gestor}): publishes and manages warehouses and the bookings made on them.
  \item \textbf{Client} (\textit{tenant}): searches, books and pays for warehouses. A client can create an organization and act as its \textbf{Organization Administrator}, or join one as a \textbf{member}.
\end{itemize}

\subsection{Sprint Plan}

Table~\ref{tab:sprint-plan} shows the Sprint Planning for the project.
Each sprint was planned and tracked in the Jira project \texttt{ES2-1-Leodega}.
The team works in two pairs, and each pair implements two sprints (Section~\ref{sec:contributions}).

\begin{table}[H]
  \centering
  \begin{tabularx}{\linewidth}{lllL{2.3cm}Y}
    \toprule
    Sprint & Start & End & Implemented by & Goal \\
    \midrule
    1 & 2026-10-01 & 2026-10-21 & Darwin Díaz, Geovanny Díaz & Shared web/mobile login (Leodeguita); register, edit, delete and moderate warehouses; catalog search and booking with simulated payment; booking management by the manager. \\
    2 & 2026-10-22 & 2026-11-13 & Darwin Díaz, Geovanny Díaz & Multi-company operation: create organizations, invite members, switch contexts and book on behalf of the organization. \\
    3 & 2026-11-23 & 2026-12-13 & Carla Gutierrez, Joao Dorado & Client--manager communication (pre-booking questions and active booking chat), incident reports and ratings. \\
    4 & 2026-12-14 & 2027-01-08 & Carla Gutierrez, Joao Dorado & External inventory integration (connection, mapping and display), Excel import and administrator metrics dashboard. \\
    \bottomrule
  \end{tabularx}
  \caption{Four-sprint plan. This report covers Sprints 1 and 2.}
  \label{tab:sprint-plan}
\end{table}

The repository history does not follow these dates exactly.
The code delivered as Sprint~1 (tag \texttt{sprint1-deploy}) is built on commit \sha{ca1799a} of 2026-10-04, and its feature commits start on 2026-09-01.
The Sprint~2 commits date from 2026-10-05 to 2026-11-04.
For this reason, this report identifies the delivered version of each sprint by its tag and commit SHA, not by calendar dates (Table~\ref{tab:deployed-versions}).
